\chapter{Sharp Hypotheses for the Four-Input RH Form}
\label{v4-arith:w3-comparison-B}

The arithmetic branch separates the following six assertions into their
established domains and their spectral or categorical hypotheses:
the P1 arithmetic-Positselski biconditional for the abelian Heisenberg
archetype, the P3 Koszul--Petersson identity, the K5
Godement--Jacquet normalisation, the Costello refinements C1--C5, the
\(\mathrm{H}\text{-}\mathrm{P}^{\mathrm{Ar}}\) comparison, and the
\(\mathrm{VB}^{*}\) implication to zero-placement.  Tate's adelic
Poisson formula gives the forward P1 direction.  The finite
Godement--Jacquet identity is a theorem on the cuspidal sph-fin core.
The Friedrichs construction gives an operator domain.  The remaining
content is precisely the passage from these domains to the full
spectral object.

\section{Six Sharp Forms}
\label{v4-arith:w3-comparison-B:overextension-list}

The six assertions have the following logical form.

\begin{table}[h]
\centering
\small
\begin{tabular}{l|p{3.3cm}|p{4.2cm}|p{4.2cm}}
No. & Assertion & Established domain & Additional hypothesis \\
\hline
1 & P1 arithmetic Positselski
(\Cref{v4-arith:thm:arithmetic-positselski},
\Cref{v4-arith:pos-costello:thm:P1-unconditional}) &
Tate-Poisson gives the forward functional-equation direction for the
abelian Heisenberg archetype. &
Completed continuous duals, derived completion, faithful lifting,
primitive zeta-character comparison, and Koszul-dual Bost--Connes. \\
2 & P3 Koszul--Petersson
(\Cref{v4-arith:kp:cor:P3-unconditional},
\Cref{v4-arith:kp:thm:main-full}) &
Archimedean Segal--Petersson identity and finite sph-fin cuspidal
Godement--Jacquet identity. &
\(\mathrm{KMS}_{1}\)/GNS identification, full Rankin--Selberg
extension, and Eisenstein summands. \\
3 & K5 absorption
(\Cref{v4-arith:bcglue-k5:thm:K5-absorption-canonical},
\Cref{v4-arith:bcglue-k5:thm:K5-adelic-consistent}) &
Cuspidal sph-fin locus after the local comparison datum. &
Compatibility with the full primary sector, including continuous
spectrum. \\
4 & Costello refinements C1--C5
(\Cref{v4-arith:pos-costello:thm:C1-closed},
\Cref{v4-arith:pos-costello:thm:C2-closed},
\Cref{v4-arith:pos-costello:thm:C3-closed},
\Cref{v4-arith:pos-costello:thm:C4-closed},
\Cref{v4-arith:pos-costello:thm:C5-closed}) &
Restricted abelian-primary, \(\mu=0\), pointed \(E_{2}\), and
\(\Lambda\)-linear domains specified below. &
The five Positselski hypotheses and the five Costello hypotheses. \\
5 & \(\mathrm{H}\text{-}\mathrm{P}^{\mathrm{Ar}}\)
(\Cref{v4-arith:hpvbconv:thm:HPAr-equivalence}) &
Friedrichs self-adjoint domain. &
Determinant divisor and Weil--Connes trace distribution for the
scaling action. \\
6 & \(\mathrm{VB}^{*}\Rightarrow\) zero-placement
(\Cref{v4-arith:rh:thm:VB-implies-F-RH},
\Cref{v4-arith:hpvbconv:thm:VBstar-converse-spectral-regularity}) &
One-sided implication after P2, P3,
\(\mathrm{H}\text{-}\mathrm{P}^{\mathrm{Ar}}\), and positive
Bost--Connes trace. &
Forward positivity; the converse is only a spectral-regularity
comparison. \\
\end{tabular}
\caption{Sharp domains for the four-input RH form.}
\label{v4-arith:w3-comparison-B:tab:overextensions}
\end{table}

The same list appears in
\Cref{v4-arith:app:C:residual-list-w2}.  The Platonic capstone is
\Cref{v4-arith:platonic:thm:RH-completion}.  Its sufficient arithmetic
path is the four-input conjunction in
\S\ref{v4-arith:w3-comparison-B:four-input-branch}.

\section{Local Forms}
\label{v4-arith:w3-comparison-B:correction-summary}

\subsection{Zero-Placement}
\label{v4-arith:w3-comparison-B:correction-23}

\begin{proposition}[one-sided zero-placement]
\ClaimStatusProvedHereConditional.  Assume P2, full P3, the
\(\mathrm{H}\text{-}\mathrm{P}^{\mathrm{Ar}}\) determinant-trace realisation,
and the positive \(\mathrm{VB}^{*}\) trace form.  Then
\Cref{v4-arith:rh:thm:VB-implies-F-RH} gives the zero-placement clause
of F-RH.  The reverse implication is not a consequence of RH; it asks
for spectral regularity of the operator model.
\end{proposition}

\begin{proof}
The implication used in
\Cref{v4-arith:rh:thm:VB-implies-F-RH} is the positive quadratic form
statement on the Bost--Connes regularised trace.  P2 and P3 identify
the test vectors and the Petersson form with the arithmetic trace
form.  \(\mathrm{H}\text{-}\mathrm{P}^{\mathrm{Ar}}\) identifies the
zero divisor and the Weil--Connes trace distribution.  Positivity then
forces the ordinates entering the Weil form to be real, hence the zeros
lie on the critical line.  None of these steps constructs the converse
positive form from zero-placement alone.
\end{proof}

\subsection{Arithmetic Positselski}
\label{v4-arith:w3-comparison-B:correction-25}

\begin{proposition}[P1 and the abelian arithmetic Heisenberg]
\ClaimStatusProvedHereConditional.  The equivalence in
\Cref{v4-arith:thm:arithmetic-positselski} and the P1 propagation in
\Cref{v4-arith:pos:thm:P1-full-unconditional} hold on the
abelian-primary Heisenberg domain after imposing the following five
hypotheses:
\begin{enumerate}
\item completed continuous duals are compatible with the relevant
co-derived and contra-derived categories;
\item projective and injective resolutions remain exact after derived
\(\Lambda\)-adic completion;
\item the Nakayama lifting used in the \(\Lambda\)-coefficient
comparison is faithful;
\item the primitive zeta-character computes the \(\kappa\)-pairing
weight;
\item the dual curved coalgebra carries the Koszul-dual Bost--Connes
interpretation.
\end{enumerate}
With these hypotheses the P1 biconditional is the arithmetic
Positselski theorem.  Without them Tate's Poisson formula supplies only
the forward implication from Koszul self-duality to the functional
equation.
\end{proposition}

\begin{proof}
Tate's formula identifies the local Fourier transforms and the global
idele quotient and gives
\(\xi(s)=\xi(1-s)\) from the abelian Heisenberg self-duality datum.
The reverse direction must lift this equality through the
\(\Lambda\)-coefficient bar--cobar adjunction.  Each of the five
hypotheses is used exactly at one such lift: topological duality,
completion, Nakayama faithfulness, the primitive trace weight, and the
Bost--Connes interpretation of the dual coalgebra.  Their conjunction
is therefore sufficient, and no smaller statement follows from Tate's
functional equation alone.
\end{proof}

\subsection{Koszul--Petersson Compatibility}
\label{v4-arith:w3-comparison-B:correction-26}

\begin{proposition}[P3 on the normalized sph-fin cuspidal core]
\ClaimStatusProvedHere.  On the normalized sph-fin cuspidal core of
\(\mathcal{V}^{\mathrm{prim}}\), the archimedean
Segal--Petersson identity and the finite Godement--Jacquet identity
assemble to the P3 Koszul--Petersson form.
\end{proposition}

\begin{proof}
The archimedean factor is the Segal Fock realisation of the Petersson
pairing.  At a finite unramified place the sph-fin vector has the
Godement--Jacquet normalisation, and the local zeta integral gives the
same Euler factor as the finite Bost--Connes trace.  The restricted
tensor product is taken over the spherical units, so the product of
the local identities is the adelic identity on the cuspidal sph-fin
subspace.
\end{proof}

\begin{remark}[full P3]
The extension from the normalized sph-fin cuspidal core to
\(\mathcal{V}^{\mathrm{prim}}\) requires the
\(\mathrm{KMS}_{1}\)/GNS identification, the Rankin--Selberg extension
of the local Petersson computation, and the treatment of Eisenstein
summands.  These three conditions are spectral hypotheses, not
formal consequences of the cuspidal calculation.
\end{remark}

\subsection{Costello Refinements C1--C5}
\label{v4-arith:w3-comparison-B:correction-34}

\begin{proposition}[Costello refinements]
\ClaimStatusProvedHereConditional.  The refinements C1--C5 imply the
P1 arithmetic Positselski biconditional on their common domain after
imposing the following five hypotheses:
\begin{enumerate}
\item C1: the \(\eta\)-algebra is obtained as a pointed colimit in the
unit-compatible \(E_{2}\) category of
\Cref{v4-arith:sw:conj:F2c-tightened};
\item C2: Parshin reciprocity holds modulo \(\Lambda\)-coboundaries,
including the \(p=2\) torsion term;
\item C3: the pro-nilpotent limit at \(\beta=1\) agrees with the
\(\mathrm{KMS}_{1}\) Tomita--Takesaki state;
\item C4: the Witt truncation has the stated classical limit and is
compatible with the archimedean Heisenberg Fock object;
\item C5: the morphism class is \(\Lambda\)-linear with trivial
\(\Gamma\)-action.
\end{enumerate}
\end{proposition}

\begin{proof}
Each C\(i\) removes one obstruction in the comparison between the
curved algebra and curved coalgebra sides.  C1 supplies the pointed
\(E_{2}\) multiplication, C2 identifies the reciprocity defect with a
\(\Lambda\)-coboundary, C3 controls the limit state, C4 fixes the
classical shadow, and C5 gives the allowed morphisms.  Their common
domain is still the domain of the arithmetic Positselski theorem, so
the five Positselski hypotheses of
\S\ref{v4-arith:w3-comparison-B:correction-25} remain present.
\end{proof}

\subsection{K5 Godement--Jacquet Normalisation}
\label{v4-arith:w3-comparison-B:correction-35}

\begin{proposition}[K5 on the cuspidal sph-fin locus]
\ClaimStatusProvedHere.  After the local comparison datum, K5 is
established on the cuspidal sph-fin locus.  Hoffstein--Lockhart
nonvanishing excludes the local zero obstruction in this domain, and
the sph-fin projection removes Eisenstein poles from the calculation.
\end{proposition}

\begin{proof}
The local comparison datum identifies the chosen sph-fin vector with
the Godement--Jacquet test vector.  The local zeta integral is then the
standard local \(L\)-factor.  On the cuspidal sph-fin locus the
Hoffstein--Lockhart nonvanishing theorem prevents the denominator
which enters the absorption normalisation from vanishing.  Since the
sph-fin projection has already removed the Eisenstein summands, no
continuous-spectrum pole enters the identity.
\end{proof}

\begin{remark}
The propagation of K5 to full Koszul--Petersson compatibility is
therefore the full P3 problem of
\S\ref{v4-arith:w3-comparison-B:correction-26}.
\end{remark}

\subsection{\(\mathrm{H}\text{-}\mathrm{P}^{\mathrm{Ar}}\) and the \(\mathrm{VB}^{*}\) Converse}
\label{v4-arith:w3-comparison-B:correction-36}

\begin{proposition}[Friedrichs domain and determinant divisor]
\ClaimStatusProvedHereConditional.  The Friedrichs construction of
\Cref{v4-arith:hpvbconv:thm:L0-Friedrichs} gives a self-adjoint
operator on its quadratic-form domain.  The Hilbert--P\'olya input is
the further assertion that the associated normal zero-divisor operator has
determinant divisor \(\mathrm{div}\,\xi_{\mathbb R}\) and
Weil--Connes trace equal to the zero distribution.  This assertion is
not a choice of zero-mode weights.
\end{proposition}

\begin{proof}
The Friedrichs theorem starts from a lower-semicontinuous positive
form and produces a self-adjoint operator.  It does not determine the
regularised determinant or the Weil--Connes trace of the scaling
action.  The Connes and Deninger forms have the required Lefschetz
shape only when the determinant divisor and the trace distribution are
the non-trivial zero divisor and distribution of \(\xi_{\mathbb R}\).
Thus the operator-domain theorem and the determinant-trace theorem are
distinct statements.
\end{proof}

\begin{proposition}[\(\mathrm{VB}^{*}\) converse]
\ClaimStatusProvedHereConditional.  The converse statement in
\Cref{v4-arith:hpvbconv:thm:VBstar-converse-spectral-regularity} is a
spectral-regularity comparison after the Fock-space identifications.
It does not imply the forward positivity statement
\(\mathrm{VB}^{*}\Rightarrow\) zero-placement.
\end{proposition}

\begin{proof}
A converse spectral-regularity theorem reconstructs regularity of the
operator model from a zero-placement input and auxiliary spectral
identifications.  The forward theorem requires positivity of the
Bost--Connes regularised trace before zero-placement is known.  The
directions therefore have different hypotheses and different content.
\end{proof}

\subsection{Primary-Source Anchors}
\label{v4-arith:w3-comparison-B:correction-A}

\begin{remark}
Li--Weil positivity, Connes' trace formula, Deninger's cohomological
form, and the Bost--Connes scaling action are four shadows of the same
problem only after the determinant divisor and trace distribution have
been specified.  The symmetrised zero measure gives the
\(\mathrm{VB}^{*}\) converse shadow; the self-adjoint scaling operator
with determinant divisor \(\mathrm{div}\,\xi_{\mathbb R}\) gives the
\(\mathrm{H}\text{-}\mathrm{P}^{\mathrm{Ar}}\) shadow.  Their
coincidence is the Hilbert--P\'olya content.
\end{remark}

\section{The Four-Input Conditional Conjunction}
\label{v4-arith:w3-comparison-B:four-input-branch}

\begin{theorem}[the four-input conditional conjunction]
\label{v4-arith:w3-comparison-B:thm:four-input-branch}
\ClaimStatusProvedHereConditional.  Assume the following four
statements.
\begin{enumerate}
\item \(\mathrm{P}_{1}^{\mathrm{ArPos}}\): arithmetic Positselski for
the abelian-primary Heisenberg domain, with the five hypotheses of
\S\ref{v4-arith:w3-comparison-B:correction-25}.
\item \(\mathrm{P}_{3}^{\mathrm{KP}}\): full-sector
Koszul--Petersson compatibility, extending the normalized sph-fin
cuspidal identity by the three spectral hypotheses of
\S\ref{v4-arith:w3-comparison-B:correction-26}.
\item \(\mathrm{HP}^{\mathrm{Ar}}\): the
\(\mathrm{H}\text{-}\mathrm{P}^{\mathrm{Ar}}\) determinant-trace realisation,
including determinant divisor \(\mathrm{div}\,\xi_{\mathbb R}\) and
the Weil--Connes zero trace distribution.
\item \(\mathrm{VB}^{*}_{+}\): the positive Bost--Connes
regularised trace form giving the forward
\(\mathrm{VB}^{*}\Rightarrow\) zero-placement implication.
\end{enumerate}
Together with the Tate-Poisson direction, P2, and Hamburger
determinacy, these four statements imply the zero-placement clause of
F-RH.  Hence the arithmetic branch implies RH under precisely this
four-input conjunction.
\end{theorem}

\begin{proof}
The arithmetic Positselski hypothesis supplies P1, because its forward
Tate direction and its reverse bar--cobar lift give the
Koszul-duality biconditional.  Full P3 identifies the adelic
Petersson form with the arithmetic trace form on the whole primary
space.  \(\mathrm{HP}^{\mathrm{Ar}}\) supplies the zero divisor and
Weil--Connes trace distribution attached to that trace.  The positive
\(\mathrm{VB}^{*}_{+}\) form is then the Weil-positive form on the zero
distribution.  By the Platonic completion theorem
\Cref{v4-arith:platonic:thm:RH-completion}, P1, P2, P3,
\(\mathrm{HP}^{\mathrm{Ar}}\), \(\mathrm{VB}^{*}_{+}\), and Hamburger
determinacy imply F-RH, hence RH.
\end{proof}

\begin{remark}[independence of directions]
\label{v4-arith:w3-comparison-B:rem:preserved}
The theorem is a sufficient implication.  It does not assert that RH
recovers arithmetic Positselski, full Koszul--Petersson compatibility,
the Hilbert--P\'olya determinant-trace realisation, or the positive
\(\mathrm{VB}^{*}\) trace form.  The
\(\mathrm{H}\text{-}\mathrm{P}^{\mathrm{Ar}}\) and
\(\mathrm{VB}^{*}\) residuals are separated by
\Cref{v4-arith:hpvbconv:prop:HPAr-det-independent}; the Positselski
and Koszul--Petersson residuals are separated by
\Cref{v4-arith:pos:residual-domain} and
\Cref{v4-arith:kp:rem:substantial-reestimate}.
\end{remark}

\begin{remark}[hypothesis support]
\label{v4-arith:w3-comparison-B:rem:changed}
The support of the four-input conjunction is local.  The P1 support is
the five arithmetic-Positselski hypotheses.  The P3 support is the
\(\mathrm{KMS}_{1}\)/GNS identification, the full Rankin--Selberg
extension, and the Eisenstein-sector calculation.  The
\(\mathrm{HP}^{\mathrm{Ar}}\) support is the determinant divisor
\(\mathrm{div}\,\xi_{\mathbb R}\) and the Weil--Connes trace
distribution.
The \(\mathrm{VB}^{*}_{+}\) support is positivity of the regularised
Bost--Connes trace.
\end{remark}

\section{Logical Form of the Essential Biconditional}
\label{v4-arith:w3-comparison-B:load-bearing}

The essential biconditional
\begin{equation*}
\kappa^{\mathrm{Ar}}_{\mathsf{G}}+(\kappa^{!})^{\mathrm{Ar}}_{\mathsf{G}}
\;=\;0\;\Longleftrightarrow\;\xi(s)=\xi(1-s)
\end{equation*}
(\Cref{v4-arith:thm:G-row-Riemann-functional-equation},
\Cref{v4-arith:thm:P1-resolution}) has asymmetric proof-theoretic
content.  The implication from arithmetic Koszul self-duality to the
functional equation is Tate's adelic Poisson formula.  The reverse
implication must reconstruct the arithmetic Positselski equivalence
for the \(\Lambda\)-coefficient abelian Heisenberg object.  It
therefore has exactly the five hypotheses of
\S\ref{v4-arith:w3-comparison-B:correction-25}.

\section{Consequences}
\label{v4-arith:w3-comparison-B:summary}

\begin{enumerate}
\item The Tate direction of P1 is unconditional; the full P1
biconditional is arithmetic Positselski.
\item P3 is established on the normalized sph-fin cuspidal core; full
P3 is the corresponding spectral extension problem.
\item K5 is established on the cuspidal sph-fin locus after the local
comparison datum.
\item C1--C5 are sufficient Costello refinements on their common
restricted domain and inherit the arithmetic-Positselski hypotheses.
\item The Friedrichs theorem gives the operator domain for
\(\mathrm{H}\text{-}\mathrm{P}^{\mathrm{Ar}}\); the determinant divisor
and trace distribution are the Hilbert--P\'olya content.
\item The \(\mathrm{VB}^{*}\) converse is a spectral-regularity
comparison and is logically distinct from the positive forward
implication to zero-placement.
\end{enumerate}
